% Gurvir Singh — one-page CV
%
% Build (from the project root, with a LaTeX distribution such as MiKTeX):
%   pdflatex -output-directory=cv/build cv/gurvir-singh-cv.tex
% then copy cv/build/gurvir-singh-cv.pdf to public/files/gurvir-singh-cv.pdf
%
% Typography follows a classic LaTeX academic CV: Computer Modern (T1 / cm-super),
% 11pt body, bold small-caps name and section headings, full-width section rules.

\documentclass[11pt]{article}

\usepackage[T1]{fontenc}
% Use the scalable Type 1 cm-super fonts (some installations do not map them by default).
\pdfmapfile{+cm-super-t1.map}
\usepackage[a4paper, left=0.8in, right=0.8in, top=0.72in, bottom=0.55in]{geometry}
\usepackage{microtype}
\usepackage{enumitem}
\usepackage{titlesec}
\usepackage[hidelinks,
  pdftitle={Gurvir Singh - Curriculum Vitae},
  pdfauthor={Gurvir Singh}]{hyperref}

\pagestyle{empty}
\setlength{\parindent}{0pt}
\setlength{\parskip}{0pt}

% ── Section headings: large bold small caps over a full-width rule ───────────
\titleformat{\section}{\Large\bfseries\scshape\raggedright}{}{0pt}{}[{\vspace{-2pt}\titlerule[1.3pt]}]
\titlespacing*{\section}{0pt}{9pt}{5pt}

% ── Entries ─────────────────────────────────────────────────────────────────
% \entry{bold left}{bold right}
\newcommand{\entry}[2]{\hspace*{6pt}\textbf{#1}\hfill\textbf{#2}\par}
% \detail{text}{right-aligned text}, indented under an entry
\newcommand{\detail}[2]{\hspace*{27pt}#1\hfill#2\par}
% \project{name}{url}{note}{stack}
\newcommand{\project}[4]{%
  \hspace*{6pt}\textbf{\href{#2}{#1}}#3\hfill\textit{#4}\par}

% Bullet lists
\setlist[itemize]{label=$\bullet$, leftmargin=34pt, labelsep=6pt,
  topsep=2pt, itemsep=1pt, parsep=0pt, partopsep=0pt}
\newenvironment{projectitems}
  {\begin{itemize}[topsep=1pt, itemsep=0.5pt]}
  {\end{itemize}\vspace{4pt}}

\begin{document}

% ── Name and contact ────────────────────────────────────────────────────────
\begin{center}
  {\LARGE\bfseries\scshape Gurvir Singh}\\[2pt]
  Third-Year Undergraduate Student
\end{center}
\vspace{-4pt}

\noindent
\hspace*{6pt}%
\begin{tabular}[t]{@{}l@{}}
  Computer Science \& Engineering\\
  (Artificial Intelligence \& Data Science)\\
  GNA University
\end{tabular}%
\hfill
\begin{tabular}[t]{@{}r@{}}
  \href{mailto:gurvir.singh.panesar@gmail.com}{gurvir.singh.panesar@gmail.com}\\
  +91 9592592881\\
  \href{https://github.com/GurvirSinghH}{github.com/GurvirSinghH}
\end{tabular}%
\hspace*{3pt}

% ── Research interests ──────────────────────────────────────────────────────
\section{Research Interests}
\begin{itemize}
  \item Machine learning and data science: data processing, analysis, and visualization
  \item Physics-informed machine learning and scientific computing (currently learning the basics)
  \item Intelligent systems: applying AI/ML methods within practical software systems
\end{itemize}

% ── Education ───────────────────────────────────────────────────────────────
\section{Education}
\entry{GNA University}{Expected 2028}
\detail{B.Tech, Computer Science \& Engineering (Artificial Intelligence \& Data Science)}{}
\detail{Third year}{CGPA: 8.13/10}
\detail{Current coursework: Data Visualization, Theory of Computation,}{}
\detail{\phantom{Current coursework: }Software Engineering, Computer Graphics}{}

% ── Selected projects ───────────────────────────────────────────────────────
\section{Selected Projects}

\project{AI Log Intelligence Platform}{https://github.com/GurvirSinghH/AI-Log-Intelligence-Platform}%
  { (ongoing)}{Python, Streamlit, Pandas, Scikit-learn, Plotly}
\begin{projectitems}
  \item Streamlit application that detects and parses Linux auth/syslog and Apache access/error logs,
        with summary statistics, keyword search, Plotly charts, and an optional Gemini incident report.
  \item Label-encoded process, module, and host fields plus PID, message length, and hour feed an
        Isolation Forest anomaly detector; K-Means then clusters the flagged entries.
\end{projectitems}

\project{Neural Forge (AI Blender Assistant)}{https://github.com/GurvirSinghH/neural-forge}%
  {}{Python, Blender Python API, LLM APIs}
\begin{projectitems}
  \item Blender add-on that turns natural-language prompts into Blender Python scripts using Gemini,
        OpenAI, local Ollama models, or any OpenAI-compatible endpoint through one API client.
  \item Optionally adds scene context (objects, active selection) to prompts, runs generated code, and
        sends runtime errors back to the model for an automatic fix.
\end{projectitems}

\project{College Knowledge Base Crawler}{https://github.com/GurvirSinghH/college-rag-assistant}%
  {}{Python, Requests, BeautifulSoup, lxml, SQLite}
\begin{projectitems}
  \item Multi-threaded crawler (worker pool and task queue) with robots.txt compliance, a shared rate
        limit, URL normalization, and retries with exponential backoff that respect Retry-After.
  \item Extracts page content and metadata to JSON, downloads linked documents, and records crawl
        state in SQLite, producing a structured dataset for a future RAG pipeline (not yet built).
\end{projectitems}

\project{CodeEditor}{https://github.com/GurvirSinghH/CodeEditor}%
  { (2nd-year project)}{Java, Swing, MySQL, JDBC}
\begin{projectitems}
  \item Desktop code editor with multi-tab editing, token-based syntax highlighting for Java, Python,
        and HTML, line numbers, and find and replace.
  \item Stores recent files, a code-snippet library, and open-tab sessions in MySQL through JDBC
        data-access objects, restoring the previous session on startup.
\end{projectitems}
\vspace{-4pt}

% ── Technical skills ────────────────────────────────────────────────────────
\section{Technical Skills}
\hspace*{6pt}%
\begin{tabular}{@{}l@{\hspace{14pt}}l@{}}
  \textit{Programming}        & Python, C, C++, Java\\
  \textit{Data and ML}        & Machine learning, data processing, data visualization\\
  \textit{Libraries and tools} & NumPy, Pandas, Scikit-learn, Plotly, Streamlit, Git, GitHub
\end{tabular}

% ── Current focus ───────────────────────────────────────────────────────────
\section{Current Focus}
\begin{itemize}
  \item Strengthening my machine learning fundamentals, to apply ML across different fields.
  \item Learning how known physics can be built into model training, e.g.\ through the loss function,
        and keeping written research notes as I study physics-informed machine learning.
  \item Early idea, not implemented: whether physics-based models combined with machine learning
        could help reason about hidden energy states in public Formula 1 telemetry.
\end{itemize}

\end{document}
